\chapter{Self-Encoding Structure}
\label{sec:self-encoding}
\label{sec:objects}
\label{sec:sandbox}
\label{sec:semantic-bootstrapping}

\section{Code at Web Scale}

A durable network must survive software change as well as storage failure. Data written today may be read by another organization years later, after schemas, libraries, runtimes, and deployment practices have moved on; preserving bytes answers only the first half of the archival question, because the reader also needs a durable account of how those bytes were structured and what operations gave them meaning. The Web already offers several approaches: JavaScript carries behavior to a widely available runtime, Protocol Buffers and similar schema-driven encodings define stable wire forms and evolution rules whose schemas supply interpretation beyond each payload, package managers and containers capture dependencies, and reproducible-build and software-transparency systems make artifacts traceable to declared sources and environments~\cite{Nikitin2017chainiac}. Each binds a different portion of the path from stored bits to behavior. At Web scale, change management crosses administrative boundaries: a recipient may accept data while declining the producer's newest software, or retain an old object long after its originating service has migrated. Reproducibility also demands explicit treatment of time, randomness, network responses, and operator choices. Self-encoding structure answers these pressures by placing a changeable definition inside the same commitment as the state it interprets, creating a historical pair that can travel, be sliced, and be cited together. The runtime becomes a smaller shared foundation, application semantics become ordinary versioned material, and transitions can expose the environmental inputs on which their results depend.

\section{Lisp as Durable Structure}

Lisp gives code and data one inspectable representation: a method definition, query, path, transition, or configuration is an expression that can be hashed, stored, transmitted, and evaluated. Homoiconicity turns semantic bootstrapping into structural composition because the network can preserve the form describing an operation with the same substrate that preserves its inputs. The implemented runtime uses s7 Scheme inside the journal; native primitives provide construction, hashing, persistence, slicing, network preparation, and controlled evaluation, while Scheme builds maps, chains, ledgers, authorization, and object protocols above them. This follows the tier-unifying intuition of single-language Web systems~\cite{Cooper2007links}, preserving explicit boundaries while reducing translations between durable state and the logic that uses it. The bootstrap terminates in native storage and evaluator semantics, which define how expressions are read, nodes are hashed, primitives are exposed, and resources are bounded. Keeping this foundation small makes it practical to inspect and reproduce while every higher layer can evolve as identified structure in journal history. A shared language still permits heterogeneous protocols: objects may implement different methods and evolve under different authorities while remaining transportable through the same node and evaluator model, with documented public behavior providing compatibility across class systems. \WIP Per-object capability assignment would extend the current shared reduced evaluation environment.

\section{Self-Encoding Objects}

A standard object is a node with executable definition on the left and durable state on the right; its digest commits to both components and to their association. Evaluating that node with \texttt{sync-eval} produces a transient callable object: a zero-argument call returns its current node, documented method dispatch exposes behavior, and object-specific paths expose state according to the object's protocol. \Cref{fig:self-encoding-transition} follows one transition as definition $D_0$ and state $S_0$ enter a controlled runtime with explicit arguments and prepared environmental values, producing a candidate with revised definition $D_1$ and successor state $S_1$. An authorized boundary chooses whether to persist that result, while the starting object remains independently identifiable and semantic change becomes historical succession.

\begin{figure}[htbp]
  \centering
  \resizebox{\textwidth}{!}{% A durable object carries its interpretation with its state.  A compatible
% controlled runtime is still required to apply that interpretation.
\begin{tikzpicture}[
  object/.style={draw=swblue, thick, rounded corners=3pt, fill=white,
                 minimum width=35mm, minimum height=31mm},
  part/.style={draw=swgray, rounded corners=1pt, minimum width=28mm,
               minimum height=8mm, align=center, font=\footnotesize},
  definition/.style={part, fill=swlightblue, draw=swblue},
  state/.style={part, fill=swlightgreen, draw=swgreen},
  runtime/.style={draw=swgold!85!black, thick, rounded corners=4pt,
                  fill=swlightgold, minimum width=38mm, minimum height=33mm,
                  text width=31mm, align=center, font=\footnotesize},
  input/.style={draw=swgray, rounded corners=2pt, fill=swlightgray,
                minimum width=28mm, minimum height=7mm,
                align=center, font=\footnotesize},
  arrow/.style={-{Latex[length=2.2mm]}, thick, draw=swgray},
  note/.style={font=\scriptsize, text=swgray, align=center},
]
  % Initial durable object.
  \node[object] (old) at (-5.0,0) {};
  \node[definition] (olddef) at (-5.0,0.55) {durable definition $D_0$};
  \node[state] (oldstate) at (-5.0,-0.55) {durable state $S_0$};
  \node[font=\small\bfseries, above=2mm of old.north] {identified object $d_0$};
  \node[note, below=2mm of old.south] {one commitment covers definition and state};

  % Controlled runtime and explicit input.
  \node[runtime] (eval) at (0,0) {\textbf{controlled runtime}\\[1mm]
                                  apply $D_0$ to the complete\\object, $a$, and $e$};
  \node[input, above=8mm of eval.north] (args) {explicit arguments $a$};
  \node[input, below=9mm of eval.south] (effects) {prepared external values $e$};
  \draw[arrow] (args) -- (eval);
  \draw[arrow] (effects) -- (eval);
  \draw[arrow] (old.east) -- node[above,note] {interpret} (eval.west);

  % Resulting durable object.
  \node[object] (new) at (5.0,0) {};
  \node[definition] (newdef) at (5.0,0.55) {resulting definition $D_1$};
  \node[state] (newstate) at (5.0,-0.55) {resulting state $S_1$};
  \node[font=\small\bfseries, above=2mm of new.north] {identified object $d_1$};
  \node[note, below=2mm of new.south] {$D_1$ may differ from $D_0$; durable only after acceptance};
  \draw[arrow] (eval.east) -- node[above,note] {result} (new.west);

\end{tikzpicture}
}
  \caption{Self-Encoding Object Transition}
  \label{fig:self-encoding-transition}
\end{figure}

Live mutation and durable mutation meet at explicit persistence: a method may update the live object's state, while the journal records the change after extracting the new node and storing it into its parent or root. Nested mutation rebuilds each containing object along the return path, keeping every predecessor available by digest and making the durable boundary visible in code. Objects compose through public protocols: deep reads traverse semantic boundaries, deep mutations rebuild ancestors, deep slicing and pruning ask each object to retain evidence appropriate to its layout, and deep merge restores compatible partial material under matching commitments. \Cref{tab:deep-operations} summarizes this implemented family. Portable definitions arrive from parties who may use different software and policy, so evaluation occurs inside a reduced environment: the s7 configuration removes broad filesystem, process, native-loading, network, random, and host-I/O authority; privileged orchestration obtains legitimate environmental values first and passes them as ordinary data; and authorization and explicit storage turn the candidate result into history.

\begin{table}[htbp]
  \centering
  \small
  \begin{tabularx}{\textwidth}{@{}>{\ttfamily\raggedright\arraybackslash}p{0.17\textwidth}>{\ttfamily\raggedright\arraybackslash}p{0.35\textwidth}X@{}}
\toprule
\normalfont\textbf{Operation} & \normalfont\textbf{Signature} & \textbf{Behavior} \\
\midrule
deep-get    & (deep-get object path) & Traverse nested public \texttt{get} methods and return the value at the path. \\
deep-set!   & (deep-set! object path value) & Set the terminal value and rebuild each containing object. \\
deep-slice! & (deep-slice! object path) & Retain evidence along the path while preserving the logical digest. \\
deep-prune! & (deep-prune! object path) & Remove evidence along the path, rebuilding parents or replacing exhausted children with stubs. \\
deep-merge! & (deep-merge! source target (path '())) & Merge digest-compatible evidence into the target or into the optional target path. \\
deep-copy!  & (deep-copy! object source-path target-path) & Read a value from one nested path and set it at another. \\
deep-call   & (deep-call object path callback) & Evaluate a callback against the selected object and return the callback result while leaving ancestors unchanged. \\
deep-call!  & (deep-call! object path callback) & Evaluate a callback, retain the mutated child node, and rebuild every ancestor. \\
\bottomrule
\end{tabularx}

  \caption{Deep Object Operations}
  \label{tab:deep-operations}
\end{table}

\section{A Controlled Semantic Web}

The project’s \texttt{standard.scm} object system demonstrates one way to grow semantics from a small evaluator: classes compile into portable nodes, methods operate through strict dispatch, trees and chains compose through public protocols, and ledgers place those objects into signed history. It provides one concrete reference architecture for portable objects, while other systems can participate by publishing durable protocols and preserving the structural invariants needed for slicing, persistence, and verification. This controlled foundation also opens two semantic surfaces: application-defined objects can carry scientific, institutional, and social structures as versioned concepts above the native runtime, while programmable queries can travel as inspectable code and execute against a fixed, authorized historical capability. Initial public queries should be read-only, metered, identified by digest, and able to return evidence for the data they touched. \WIP Installation, invocation, proof round trips, authorization, metering, and migration remain the acceptance tests for both facilities. Together they complete the progression from integrity to meaning, which Chapter~5 projects outward as a Web whose coordinates, identity, and trust begin from a selected journal. At network scale, control concentrates in three areas:

\begin{itemize}
  \item \textbf{Deterministic transition.} Successors derive from the starting object, definitions, explicit arguments, evaluator semantics, and controlled primitives; time, randomness, remote responses, and operator choices become prepared inputs. \WIP The current system-time procedures require removal, mediation, or recording.
  \item \textbf{Resource governance.} Evaluation needs bounds on time, allocation, recursion, reads, result size, and durable mutation, with atomic failure at a bound~\cite{Levy1995safe,Necula1997proof}. \WIP Object-specific budgets and capabilities would extend the current shared reduced environment.
  \item \textbf{Explicit evolution.} Public methods, return shapes, sentinels, and layouts form compatibility contracts, and migrations record the relation between old and successor definitions and state. \WIP Complete migration must cover authority, rollback, nested partial objects, validation, and long-term runtime compatibility.
\end{itemize}